% Generated draft: requires mathematical and source review.
\documentclass[11pt]{article}
\usepackage[margin=25mm]{geometry}
\usepackage{fontspec}
\setmainfont{Latin Modern Roman}
\usepackage{amsmath,amssymb,mathtools,unicode-math}
\setmathfont{Latin Modern Math}
\usepackage{enumitem,needspace,xurl,hyperref}
\setlength{\parindent}{0pt}
\setlength{\parskip}{5pt}
\setlength{\emergencystretch}{3em}
\newcommand{\N}{\mathbb{N}}
\newcommand{\Z}{\mathbb{Z}}
\newcommand{\Q}{\mathbb{Q}}
\newcommand{\R}{\mathbb{R}}
\newcommand{\C}{\mathbb{C}}
\newcommand{\F}{\mathbb{F}}
\newcommand{\A}{\mathbb{A}}
\newcommand{\PP}{\mathbb{P}}
\newcommand{\E}{\mathbb{E}}
\newcommand{\eps}{\varepsilon}
\newcommand{\dd}{\,\mathrm{d}}
\begin{document}
\section*{UnsolvedMath ID: 20001905}
\textbf{Source problem code:} catalogue entry
\subsection{Definitions and mathematical statement}
Question 7.2.1 asks whether homotopic closed geodesics in a manifold without conjugate points have the same length. Under the complete Riemannian interpretation of the aid, let \(c_1,c_2\) be closed geodesics in the same free homotopy class and let \(L_i\) be their lengths, including their traversing multiplicities. The reported affirmative statement is
\[
 L_1=L_2.
\]
The aid cites Climenhaga--Knieper--War, Lemma 2.29, \emph{Advances in Mathematics} 376 (2021). Its proof description synchronizes lifts by a deck transformation and uses globally minimizing geodesics in the universal cover. Completeness is an explicit hypothesis of this interpretation, not present in the short original question.

The aid's further quantitative statement tracks powers. Let \(m,n\geq1\), and suppose \(c_1^m\) and \(c_2^n\) are freely homotopic, possibly after reversing one orientation. Choose oriented lifts whose periods are synchronized by one deck isometry \(\gamma\), and points \(x_1,x_2\) on them. Define, for integers \(k\geq1\),
\[
 E_1(k)=kmL_1-d(x_1,\gamma^k x_1),\qquad
 E_2(k)=knL_2-d(x_2,\gamma^k x_2).
\]
The reported finite-iteration inequality is
\[
 \left|k(mL_1-nL_2)-(E_1(k)-E_2(k))\right|
 \leq2d(x_1,x_2).
\]
Thus sublinear defects \(E_i(k)=o(k)\) force \(mL_1=nL_2\), not necessarily equality of the primitive lengths when the powers differ. With stable translation length
\[
 \tau(\gamma)=\lim_{k\to\infty}\frac{d(x,\gamma^k x)}k,
\]
the limiting defects are \(mL_1-\tau(\gamma)\) and \(nL_2-\tau(\gamma)\). Their difference is exactly the corresponding period-length difference. In the no-conjugate-points setting, global minimizing of the synchronized lifted geodesics removes the defects and recovers the stated rigidity. The quantitative defect language is not a claim that arbitrary closed geodesics in arbitrary metrics minimize globally.

The record then transitions to a separate large-scale parallel-geodesic discussion. It expresses hope for a flat-plane theorem in a no-conjugate-points manifold with a sufficiently large family of ``parallel'' geodesics, potentially yielding a totally geodesic isometric immersion of the Euclidean plane. No precise family hypothesis or theorem is supplied. It asks in particular whether a complete Riemannian metric on \(\mathbb R^2\) satisfying Euclid's fifth postulate must be flat: for every complete geodesic image and every point outside it, exactly one nontrivial unoriented geodesic image through the point is disjoint from the first. The source says this postulate implies no conjugate points. The preceding theorem about closed geodesic lengths does not settle this separate plane-rigidity question, and the record's solved label is not applied to it.
\subsection{Short English statement}
The supplied literature reference gives equal lengths for freely homotopic closed geodesics in the complete no-conjugate-points setting. The aid also records a power-sensitive minimizing-defect estimate. Its neighboring discussion of parallel geodesics and Euclid's fifth postulate is retained as a separate, incompletely specified rigidity question.
\subsection{Sources}
\begin{itemize}
\item Local source record: catalogue entry, catalogue entry; canonical \texttt{aim-geometry-notes.json}, \texttt{notes[242]}, source item 7.2, extracted Question 7.2.1. Title: ``Length rigidity in a free homotopy class.''
\item Keith Burns and Vladimir S. Matveev's AIM problem list, workshop ``Geodesics,'' supplied reference \url{https://aimath.org/pastworkshops/geodesicsproblems.pdf}. The verification identifies the question followed by a transition paragraph. This draft uses the supplied text without accessing a PDF.
\item Supplied literature citation: Climenhaga--Knieper--War, Lemma 2.29, \emph{Advances in Mathematics} 376 (2021). A paper title and identifier are not supplied. The quantitative defect extension belongs to research attempt 1.
\item Registry dated 14 August 2026: ``solved,'' ``SOLVED-IN-LITERATURE,'' statement status ``exact''; attempt full solution, research-run status solved, rigor ``proved,'' novelty confidence low. These concern the main length question in the report. The machine-generated, non-peer-reviewed aid requires independent verification before citation.
\end{itemize}
\end{document}
